% =====================================================================
% STAGED ARTIFACT 3/3 -- plus-page camera-ready (+1 page, limit 8, body 7)
% response_to_reviewers.tex
% Purpose: the compact "Response to reviewer comments" block that is the
%          STATED reason the workshop grants the +1 page. Maps each of the
%          three accept-decision workshop reviewers' main criticisms to the
%          exact camera-ready change that addresses it, citing our own
%          sections. Honest and concrete; flags what is in-progress vs. done.
%
% SOURCE of the six criticisms (per task brief; the accept-decision reviewers
% raised): RVR triviality; small samples; quantization confound (all-4bit);
% no downstream task-success / latency; BFCL-only (single benchmark);
% contamination. Grouped below under three reviewers (R1/R2/R3). If the real
% reviews assign these differently across reviewers, re-bucket the rows -- the
% mapping criticism->fix is what matters and is reviewer-agnostic.
%
% HONESTY FLAGS (do not soften beyond what the data supports):
%   - bf16/8-bit quant control is IN PROGRESS, not landed. The appendix
%     NUMBERS_TODO (A1_appendix.tex ~L153) lists Qwen3-8B-8bit / -bf16 /
%     14B-8bit as still running. Phrase as "in progress," not "we show."
%   - Downstream: we report strict-pass deltas (Sonnet -11.7pp, Haiku -6.7pp;
%     05_results.tex L200-207), abstention-by-rerouting (trace, L243-251), and
%     RVR round-trip / token budget (L249-251). We do NOT log wall-clock
%     latency separately -- say "provider round-trips and token overhead,"
%     NOT "latency in ms."
%   - Second benchmark (tau-bench): plumbing verified (0/24 pooled, smoke
%     only -- see _staging/INTEGRATION_NOTES.md L202-204); NO paper-ready
%     pass-rate. Cite as "wired and smoke-checked, results pending," do NOT
%     claim a second non-zero site.
%   - Contamination: still acknowledged as unresolved (07_limitations.tex L5);
%     we report it as a bounded caveat, not a closed gap.
%
% INTEGRATION:
%   Place this block LAST in the body, after Section~\ref{sec:limitations}
%   and before the LLM-disclosure section, as an unnumbered section so it
%   reads as the camera-ready addendum the +1 page is granted for. In main.tex,
%   add  \input{sections/_response_to_reviewers}  between \section{Discussion and Limitations}
\label{sec:limitations}
\label{sec:limitations:rvr-null}

\input{sections/_inbody_deployer}

Table~\ref{tab:deployer} condenses the tier-conditional recommendation
into one operator-facing artifact; the rest of this section states what
the audit does and does not license.

We cannot fully disentangle the Anthropic--Qwen3 gap: parameter count, training recipe, and provider guardrails co-vary, and commercial-vs-open still bundles the serving path with the weights. Quantization, however, is no longer a confound: a precision ladder (4-bit/8-bit/bf16) on Qwen3-8B shows fabrication persists at full precision (Section~\ref{sec:results:controls}), so it is not a quantization artefact, and a second open lineage (Llama-3.1-8B) reproduces the gap. The Anthropic zero-event regime is indistinguishable from BFCL v4 contamination, since v4 predates every model tested and we ran no paraphrase or membership-inference control. RVR's $20\!\to\!0$ prevention is invariant to quantization, as $\mathrm{C}_0$ and $\mathrm{C}_1$ share weights. On Anthropic, RVR is a call-path no-op yet strict-pass drops (Section~\ref{sec:results:rvr}), and our Newcombe-hybrid TOST is underpowered at $N\!=\!60$. We recommend $\mathrm{C}_{0.7}$/RVR only on tiers with audited $\mathrm{TFR}\!>\!0$, paired with executor-side allow-listing. Scope: the probe tests registry-shape, not runtime-shape (Czapla's \texttt{read\_secret()} collision needs ambient-reachability BFCL lacks); TFR ignores latency, abstention, and pass-rate; and we claim no scaling law. We do show the fabrication gap in a second open family (Llama-3.1-8B) and that RVR sharply reduces it there (to $0.41\%$, with both residual calls caught by the rejection layer rather than reaching the executor), but with only two open lineages, and a cross-family null looser than the per-tier-matched Qwen3 sweep, we do not claim general cross-family transfer. A second benchmark bounds the finding's reach rather than extending it. On $\tau$-bench---multi-turn, with a simulated user---the commercial null holds ($0/2{,}244$ pooled across gpt-4.1, Sonnet~4.6, and Haiku~4.5, $\leq\!0.16\%$), but our target-removed probe did not elicit open-weight fabrication: Qwen3-8B and 14B emitted $191$ valid in-registry calls with zero fabrications ($\leq\!1.91\%$) while completing none of the $24$ probed $\tau$ tasks. We read this as benchmark-structure dependence---$\tau$'s user-sim and larger task registries do not reproduce BFCL's clean ``expected tool absent'' pressure---and therefore scope the open-weight gap to BFCL-style target-removed registries, not to tool use in general. Because those $\tau$ tasks also went uncompleted, the null may reflect too little successful tool engagement to fabricate on rather than genuine robustness; we read $\tau$ as a scope boundary, not a negative result. Breadth is bounded by tooling for the same reason: four further open families (Mistral-7B-v0.3, Gemma-2-9B, Phi-3.5-mini, DeepSeek-R1-Distill-7B) emitted no parseable tool calls under our text-parse harness and yield no rate, so the open-weight evidence rests on the three lineages above (Qwen3, Qwen2.5, Llama).

%   and \section*{LLM/Agent Usage Disclosure}
\label{sec:llm-disclosure}

\paragraph{Experimental subjects.} The paper evaluates eleven LLMs as
subjects of measurement: five Anthropic 4.x versions (Opus~4.7,
Sonnet 4/4.5/4.6, Haiku~4.5) via Anthropic's tool-use API, and six
Qwen3-Instruct sizes (0.6B--32B) as 4-bit MLX builds on a $32$~GB
Apple Silicon machine. All are treated as black-box targets;
version snapshots, dated aliases, and decoding configurations are
pinned in the supplementary manifest.

\paragraph{Drafting, code, and statistics.} The paper text and the
harness (adapters, intervention conditions, statistical scripts,
trace logger, cost meter) were drafted with LLM-based assistance
under author supervision. All claims, numerical results,
methodological decisions, and statistical tests (Fisher's exact,
Newcombe-hybrid TOST, Holm--Bonferroni, BCa cluster-bootstrap,
Friedman+Nemenyi) are reproducible from the released harness and
its $144$ unit tests; no $p$-value appears without a code path that
regenerates it. Contributions, experimental design, scope claims,
and limitations were authored under direct human supervision.
. Uses only \toprule/\midrule (booktabs) and
%   standard ICML packages -- no new preamble. Single column.
% =====================================================================

\section*{Response to Reviewer Comments}
\label{sec:response}

We thank the reviewers. This camera-ready page records, point by point,
how each main criticism is addressed; we flag in-progress items honestly
rather than overclaim.

\paragraph{R1: RVR is trivial / a membership check is obvious.}
We agree the mechanism is simple, and we now frame that as the point
rather than hiding it. RVR is the message-level analog of constrained
decoding for the closed-API regime where logits are unavailable
(Section~\ref{sec:method:rvr}): an $O(1)$ in-process membership check
plus one structured re-prompt, training-free, composable with RL methods.
Its value is empirical and bounded by measurement, not by novelty---the
per-call TFR guardrails (system-failure exclusion, per-call rather than
per-task denominators, cluster-bootstrap intervals;
Section~\ref{sec:method:tfr}) are what let a $14\!\to\!0$ prevention be
read as signal rather than trace-length noise. The four-condition ablation
(Section~\ref{sec:results:rvr}, Figure~\ref{fig:ablation}) further shows
the work is done by the structured-error envelope $\mathrm{C}_{0.7}$
alone; echoing the full registry ($\mathrm{C}_1$) adds nothing detectable
(Newcombe $95\%$ CI on $\mathrm{C}_1\!-\!\mathrm{C}_{0.7}$:
$[-1.5,+1.4]$\,pp), so we report a deployable, name-leak-free variant, not
a heavyweight method.

\paragraph{R1: Small samples / underpowered.}
We have stopped leaning on point estimates and now report exact bounds
throughout. Every zero-event cell carries a Clopper--Pearson $95\%$ upper
bound (pooled commercial frontier $0/2{,}117$ at $\leq\!0.14\%$ and
$0/2{,}578$ at $\leq\!0.12\%$; Section~\ref{sec:results:probe},
Table~\ref{tab:frontier}). The RVR contrast is tested with Fisher's exact
on the pooled $14/973$ vs.\ $0/945$ table ($p\!=\!7.1\!\times\!10^{-5}$,
surviving Holm--Bonferroni) and we state plainly that each per-tier test
is individually underpowered (Section~\ref{sec:results:rvr}). The
Anthropic strict-pass non-inferiority test is reported as
\emph{undetermined} at $N\!=\!60$ ($1\!-\!\beta\!<\!0.30$), not as a pass.
An $N\!\geq\!100$ event-multiplication run on the hot $8$B/$14$B cells is
in progress to move the per-tier tests off the floor.

\paragraph{R2: Quantization confound---everything is $4$-bit.}
This is the sharpest objection and we treat it as open. The body no longer
asserts the curve is ``a capability effect, not a compression artifact'';
Section~\ref{sec:limitations} states we cannot exclude a
quantization-damage account, and the revised scaling caption
(Figure~\ref{fig:scaling}) reads the shape descriptively only. The
deciding control---a bf16 (and $8$-bit) Qwen3-$14$B/$8$B arm at matched
$N$, with the per-checkpoint quantization recipe disclosed---is
\emph{in progress} (Appendix, \textsc{numbers\_todo}); we will report
whether the $1.87\%$ peak survives or collapses. We also note the one
claim that is confound-\emph{free} regardless of the control: RVR's
$14\!\to\!0$ prevention holds weights fixed between $\mathrm{C}_0$ and
$\mathrm{C}_1$, so it cannot be a quantization artifact
(Section~\ref{sec:limitations}).

\paragraph{R2: No downstream task-success or latency.}
We now report the operational deltas RVR induces, not only fabrication
counts. Strict task-pass under $\mathrm{C}_1$ on the Anthropic frontier
drops Sonnet $60/60\!\to\!53/60$ ($-11.7$\,pp) and Haiku
$60/60\!\to\!56/60$ ($-6.7$\,pp) (Section~\ref{sec:results:rvr}); we are
explicit that at $N\!=\!60$ this drop is not distinguishable from noise.
For the abstention/behavioral channel, the trace analysis
(Section~\ref{sec:results:trace-inspection}) shows agents \emph{reroute}
through in-registry tools rather than fabricate (e.g.\ a five-call
in-registry chain in place of a fabricated \texttt{mv\_recursive()}). For
cost, RVR adds zero provider round-trips on hits and exactly one on
misses, with a rendered registry of $\sim\!100$--$200$ tokens
($<\!5\%$ of context) at $\leq\!25$ tools
(Section~\ref{sec:results:trace-inspection}). We do not log wall-clock
latency in milliseconds and do not claim a latency number; the cost is
reported as round-trips and token overhead.

\paragraph{R3: Single benchmark (BFCL-only).}
We acknowledge the breadth limit and have built toward a second site
rather than asserting generality. A second harness ($\tau$-bench) is wired
end-to-end and smoke-checked through the same TFR pipeline; paper-ready
pass-rate and TFR numbers on it are pending and we make no
second-benchmark claim until they land. Until then we scope every positive
result to BFCL v4 multi-turn explicitly (Section~\ref{sec:limitations})
and frame the cross-vendor split as a serving-and-model bundle on one
harness (Table~\ref{tab:frontier}), not a benchmark-independent property.

\paragraph{R3: Contamination.}
We do not claim to have ruled this out. BFCL v4 predates every model
tested, so the Anthropic zero-event regime is, as we state, observationally
indistinguishable from training-set contamination; we ran no paraphrase or
membership-inference control (Section~\ref{sec:limitations}). We therefore
present the commercial null as an \emph{upper bound under possible
contamination}, and we note that the open-weight non-zero band and the
RVR prevention result do not depend on the contamination question---a
contaminated model that still fabricates ($6.10\%$ on Qwen2.5-7B;
Section~\ref{sec:results:scaling}) shows the probe is not merely re-reading
memorized registries.
